\documentclass[11pt]{article}
\usepackage[margin=2.4cm]{geometry}
\usepackage{fontspec}
\setmainfont{DejaVu Sans}
\usepackage{juliamono}
\usepackage{fvextra}
\usepackage{xcolor}
\usepackage{xurl}
\usepackage{hyperref}
\usepackage{amssymb}
\newcommand{\checkbox}{$\square$}
\newcommand{\answerchoices}{\checkbox\ YES \quad \checkbox\ NO \quad \checkbox\ CANNOT DETERMINE}
\newcommand{\decisionexplanation}{%
  \par\smallskip\textbf{If NO, give the reason for the mismatch (required). If
  CANNOT DETERMINE, state what context or evidence is missing:}\par
  \rule{\linewidth}{0.4pt}\par\vspace{7pt}\rule{\linewidth}{0.4pt}\par}
\setlength{\parindent}{0pt}\setlength{\parskip}{5pt}
\setlength{\emergencystretch}{2em}
\DefineVerbatimEnvironment{Vsrc}{Verbatim}{fontsize=\small,frame=single,
  breaklines=true,breakanywhere=true}
\DefineVerbatimEnvironment{Vlean}{Verbatim}{fontsize=\scriptsize,frame=single,
  numbers=left,numbersep=5pt,breaklines=true,breakanywhere=true}
\hypersetup{colorlinks=true,urlcolor=blue}

\begin{document}
\begin{center}{\LARGE Proof Verification Document}\\[2pt]{\large SE\_cs\_81391\_\_q}\\[2pt]{\small Stack Exchange, question arm (confirmed)}\end{center}
\begin{center}\fcolorbox{blue!60!black}{blue!4}{\parbox{0.92\linewidth}{\textbf{Make three independent decisions.} Follow the evidence chain: original source $\longrightarrow$ challenged claim $\longrightarrow$ proposed counterexample $\longrightarrow$ Lean theorem. A record is valid only when all three links are correct. Complete the human sign-off before reading the automated review appendix; those earlier results are background, not an answer key.}}\end{center}
\textbf{Use only this document, except for the public source link.} Everything needed for Questions 2 and 3 is included here. For Question 1, open the public source link when one is shown. If there is no public link, the complete stored problem and solution are included here. If the information is still not enough, choose CANNOT DETERMINE.
For each question, choose \textbf{YES}, \textbf{NO}, or \textbf{CANNOT DETERMINE}. Choose CANNOT DETERMINE if the document does not provide enough information for you to decide; do not mark YES if you are uncertain. If you find a specific mismatch, choose NO and explain it in the space below that question. Do not rewrite the claim, add assumptions, or substitute a different counterexample.
\textbf{How copied text is shown.} The problem, quoted passage, and proposed counterexample come from the stored record. The Lean theorem statement and Lean files are copied exactly from the certificate. Some long lines continue on the next printed line; this is only page formatting and does not change the copied text. Line numbers are added only to help you find passages. If the stored record contains a character that cannot be printed, the document shows \texttt{[U+XXXX CONTROL]} in its place.
\begin{description}
\item[Challenged claim.] The exact statement from the source that this record says is false.
\item[Proposed counterexample.] A concrete case that should satisfy the claim's assumptions but make its conclusion false.
\item[Lean certificate.] A Lean file proving a formal theorem. Successful compilation proves the theorem as written, not that it is the right translation of the source.
\end{description}
\section*{Question 1 --- Does this proof verification document accurately represent the source?}
\textbf{What to decide.} Compare the challenged passage below with the original source---either the public link or the complete source text in this document. Answer YES only if they make the same claim and include the same assumptions or conditions. If this document changes the meaning, leaves out a condition, or points to a different passage, answer NO.
\subsection*{Original source}
\textbf{Source record:}
\begin{Vsrc}
qedbench_certificate_availability_triage/data/records_se.jsonl
\end{Vsrc}
\textbf{Public source:}\par
\url{https://cs.stackexchange.com/questions/81391/proof-if-the-list-is-then-k-sorted-for-some-smaller-integer-k-then-the-list-r}\par
\subsection*{Problem or question from the source (verbatim)}
\begin{Vsrc}
Proof: "If the list is then k-sorted for some smaller integer k, then the list remains h-sorted"
\end{Vsrc}
\subsection*{Challenged passage recorded from the source (stored text)}
\begin{Vsrc}
Say that after first pass $x_0$ switches place with $x_m$, so the final array $B$ starts with the entry $x_m$ regardless of the position of $x_0$ (which can further switch place with $x_{2m}$ and so on). Thus, $x_m < x_0$.
\end{Vsrc}
\subsection*{Surrounding source lines (stored text with added line labels)}
The labels refer to lines in the solution text stored with the source record; the source-line content is not transliterated.
\begin{Vsrc}
0012 | > sorted list in the end.
0013 | 
0014 | [[Source](https://en.wikipedia.org/wiki/Shellsort)]
0015 | 
0016 | *What is the mathematical proof for the quoted claim (in bold)?*
0017 | 
0018 | **An attempt at the proof** (for the case when number of elements is $2n+1$):
0019 | 
0020 | 
0021 | ----------
0022 | 
0023 | 
0024 | My initial array is $A = \{x_0, \cdots, x_n\}$ and i have $x_0 < x_n$. Say that after first pass $x_0$ switches place with $x_m$, so the final array $B$ starts with the entry $x_m$ regardless of the position of $x_0$ (which can further switch place with $x_{2m}$ and so on). Thus, $x_m < x_0$. Now, in the final pass $x_n$ can switch place with $x_k$ (which, before the last pass, is in the $(n-m)$-th position in the array). We don't care about $k$, but we just know that $x_n < x_k$. So the last entry in the final $m$-sorted array $B$ is $x_k$.
0025 | Sum it up: $x_m < x_0 < x_n < x_k$.
0026 | 
0027 | That means $B$ is $n$-sorted.
0028 | If we assumed the switchings of $x_0$ and $x_n$ didn't happen during $m$-sorting, we are trivially done, because the starting and final entries of $B$ are $x_0$ and $x_n$ respectively, and since $x_0 < x_n$, $B$ is $n$-sorted.
0029 | 
0030 | Say $A = \{x_0, \cdots, x_n, \cdots, x_{2n}\}$ is $n$-sorted. We can assume the subarray $\{x_0, \cdots, x_n\}$ is $m$-sorted WLOG and because of what we proved earlier. Suppose now that while $m$-sorting, $x_n$ and $x_{n+m}$ are switched; then $x_{n+m} < x_n$. And then on the last entry, $x_k$ (in the $2n-m$ th position) and $x_{2n}$ are switched. Then $x_{2n} < x_k$.
0031 | 
0032 | After the $m$-sorting is over, the $n$-gap subarray starting at $x_0$ is $\{x_0, x_{n+m}, x_k\}$. We have $x_0 < x_m < x_{n+m} < x_n < x_{2n} < x_k$, where $x_m < x_{n+m}$ is because $A$ is $n$-sorted.
0033 | 
0034 | So, this is proof that if a $2n+1$ array is $m$-sorted after being $n$-sorted, it remains $n$ sorted.
0035 | 
0036 | 
\end{Vsrc}
\textbf{A YES means:} the proof verification document faithfully represents both the source context and the challenged claim, without dropping a condition or changing its meaning.
\section*{Question 2 --- Does the proposed counterexample disprove that claim?}
\textbf{Claim to test (repeated for convenience).}
\begin{Vsrc}
Say that after first pass $x_0$ switches place with $x_m$, so the final array $B$ starts with the entry $x_m$ regardless of the position of $x_0$ (which can further switch place with $x_{2m}$ and so on). Thus, $x_m < x_0$.
\end{Vsrc}
\subsection*{How the record says the argument uses this claim}
This is the record's description, not an established fact; check it against the source.
\begin{Vsrc}
The proof asserts that if x_0 switches place with x_m during the first pass of m-sorting, then the final m-sorted array B starts with the entry x_m regardless of where x_0 moves afterwards. This is used to conclude the chain x_m &lt; x_0 &lt; x_n &lt; x_k and hence that B is n-sorted. In fact B's first entry is the minimum of the entire gap-m subsequence through position 0, which can be x_{2m} (or a later element), not x_m.
\end{Vsrc}
\subsection*{Proposed counterexample (exact values or construction)}
\begin{Vsrc}
Non-degenerate witness (genuine gap m >= 2), smallest of its kind found by exhaustive enumeration:

  N = 5, m = 2, A = (x_0,...,x_4) = [3, 4, 2, 5, 1]

  Gap-m chain through position 0: x_0, x_2, x_4 = 3, 2, 1.
  Hypothesis of the claim holds: at the first insertion step (i = m = 2) we compare x_2 = 2 with x_0 = 3; since 2 < 3 they exchange places, i.e. "x_0 switches place with x_m during the first pass".
  After the full m-sort (gap-2 insertion sort, equivalently sorting each residue chain in place):

  B = m-sort(A) = [1, 4, 2, 5, 3]

  The claim says B[0] == x_m = 2; here B[0] = 1 != 2. In fact B[0] = 1 = x_4 = x_{2m}, the minimum of the whole gap-m chain {x_0, x_m, x_{2m}} = {3, 2, 1} — exactly the failure mode described in the target claim.

Smallest witness overall (degenerate gap m = 1, i.e. plain insertion sort): N = 3, A = [3, 2, 1]; x_1 = 2 < 3 = x_0 so they swap on the first step, but B = [1, 2, 3] and B[0] = 1 = x_2 != x_1 = 2.

The script also verified, over all 36243 counterexamples found, that B[0] always equals min{x_0, x_m, x_{2m}, ...} (0 exceptions), confirming the stated diagnosis rather than merely the negation.
\end{Vsrc}
\subsection*{Why the record says the counterexample works}
Check this explanation independently; it is evidence to assess, not an answer key.
\begin{Vsrc}
The asserted implication is: if during the first insertion step of m-sorting x_0 and x_m exchange places (equivalently x_m < x_0), then the final m-sorted array B satisfies B[0] = x_m.

This is false. m-sorting sorts each gap-m residue chain, so B[0] = min{x_0, x_m, x_{2m}, x_{3m}, ...} — the minimum over the entire chain through position 0, not merely the smaller of x_0 and x_m. Whenever some later chain element is smaller than x_m, the element x_m is itself displaced further right by a subsequent insertion, and it is that later element that ends up in position 0. The parenthetical in the proof ("which can further switch place with x_{2m} and so on") tracks x_0 moving right but overlooks that x_m moves right in exactly the same way.

Witness: A = [3, 4, 2, 5, 1] with m = 2. Hypothesis holds (x_m = 2 < 3 = x_0, they swap at i = m). Recomputed: B = [1, 4, 2, 5, 3], so B[0] = 1 whereas the claim asserts B[0] = x_m = 2. Claim says B[0] = x_m = 2; here B[0] = 1 != 2, and B[0] = x_{2m} = x_4 = 1.

Consequence for the proof: the identification of B's first entry with x_m is unjustified, so the chain of inequalities x_m < x_0 < x_n < x_k that is derived from "B starts with x_m" is not supported by this argument. (The weaker fact x_m < x_0 does follow directly from the swap hypothesis, but the step that names B[0] as x_m — and hence the link into the rest of the chain — fails.)

Exhaustiveness note: the counterexample is verified by exact integer comparisons only; the m-sort was computed by two independent implementations (gap-m insertion pass and per-residue-chain sorting) which were asserted to agree on every input examined.
\end{Vsrc}
\subsection*{How to decide}
\begin{enumerate}
\item List every assumption that applies to the challenged claim.
\item Check that the proposed example satisfies every one of those assumptions.
\item Recompute the claimed outcome and the actual outcome using exact reasoning.
\item Confirm that the conflict is with the challenged claim itself, not an unrelated sentence.
\end{enumerate}
\textbf{A YES means:} the example is fully specified, satisfies all relevant assumptions, and makes the challenged claim false. If a necessary value, condition, or calculation is missing, choose CANNOT DETERMINE.
\section*{Question 3 --- Does the final Lean theorem statement verify this same counterexample?}
\textbf{What to decide.} A Lean file can compile while proving the wrong theorem. Compare these three links independently:
\begin{enumerate}
\item The Lean definitions represent the same mathematical objects, assumptions, and operations as the challenged claim.
\item The Lean construction uses the same concrete values or example shown in Question 2.
\item The final theorem statement asserts that the challenged claim fails on that example, not merely a weaker or nearby statement.
\end{enumerate}
\subsection*{Selected final Lean theorem statement (verbatim UTF-8 excerpt)}
This is an exact contiguous slice of the certificate file, from the theorem keyword through the original \texttt{:=} token where its proof begins. It is shown first so the formal target is not buried in the complete file.
\begin{Vlean}
theorem claim_false_on_witness : claimAt mW A = false :=
\end{Vlean}
\subsection*{Project-local imported Lean files (complete, verbatim UTF-8)}
This section follows the certificate's import commands recursively and prints every repository-owned Lean source file on which the final theorem statement depends. Each file is copied in full, including its definitions and controls. External Mathlib, Std, and Lean Init source is not duplicated.
No project-local imported Lean files were found. External Mathlib, Std, and Lean Init source is not duplicated in this proof verification document.
\subsection*{Complete counterexample Lean file (verbatim UTF-8)}
The complete Lean file is included below so you can check every definition and supporting result used by the final theorem statement. Use it together with the project-local imported Lean files above.
\textbf{Lean certificate file:}
\begin{Vsrc}
partB_feasibility/se_q_certificates/SE_cs_81391__q.lean
\end{Vsrc}
\begin{Vlean}
/-!
# SE_cs_81391__q -- self-contained Lean 4 refutation certificate

RECORD ID : `SE_cs_81391__q`
QUESTION  : https://cs.stackexchange.com/questions/81391
            (URL inferred mechanically from the record id `SE_cs_<id>__q`; the record itself
            carries no URL field.  The flagged proof lives in the QUESTION body -- this is a
            `__q` record -- so what is refuted below is the asker's own argument.)
ROUTE     : SELF_CONTAINED_LEAN (no imports at all; Lean core `Init` only; kernel `decide`
            throughout -- no `native_decide`, no `partial`, no `sorry`, no `axiom`,
            no well-founded recursion)

## Route note (why not the V2 claim DSL)

`qedbench_certificate_availability_triage/lean_certificates/QEDTemplateV2.lean` was checked
constructor-by-constructor.  `WitnessV2` carries graphs, a weighted digraph, `Nat` lists and
pair lists, and the only sorting-flavoured constructor is `quicksortSortedNLogN (n c : Nat)`,
whose semantics (`quicksortSortedNLogNB`) is a COMPARISON-COUNT bound for first-pivot
quicksort on the already-sorted input `[0 .. n-1]`; it takes no witness array at all.  There
is no constructor mentioning a gap/increment, no constructor whose hypothesis is "positions
0 and m exchange during the first insertion pass", and no constructor whose conclusion reads
off an entry of an algorithm's OUTPUT array.  Neither hypothesis nor conclusion of the claim
under test is expressible over the existing constructors, so the route is
SELF_CONTAINED_LEAN rather than V2_DSL.  (Nothing in the template is modified or imported.)

## The claim being refuted

The record supplies no verbatim source sentence; what is formalised here is the record's own
precise rendering of the flagged inference, quoted verbatim from `refutation_explanation`:

    "The asserted implication is: if during the first insertion step of m-sorting x_0 and
     x_m exchange places (equivalently x_m < x_0), then the final m-sorted array B satisfies
     B[0] = x_m."

Here "m-sorting" is the gap-`m` (increment-`m`) insertion pass of Shellsort: for
`i = m, m+1, ..., N-1`, insert `x_i` backwards along its residue chain `i, i-m, i-2m, ...`.
The record also states the flaw and the true value, again verbatim:

    "m-sorting sorts each gap-m residue chain, so B[0] = min{x_0, x_m, x_{2m}, x_{3m}, ...}
     -- the minimum over the entire chain through position 0, not merely the smaller of x_0
     and x_m.  Whenever some later chain element is smaller than x_m, the element x_m is
     itself displaced further right by a subsequent insertion, and it is that later element
     that ends up in position 0.  The parenthetical in the proof ('which can further switch
     place with x_{2m} and so on') tracks x_0 moving right but overlooks that x_m moves
     right in exactly the same way."

Both halves are certified below: the claim's conclusion is refuted (`claim_false_on_witness`)
AND the record's diagnosis `B[0] = min(chain through 0)` is confirmed, on the witness and
across every exhaustive sweep in this file (`witness_diagnosis`, `nc_repaired_*`).

## Witness (the record's `witness_description`, transcribed unchanged)

    N = 5,  m = 2,  A = (x_0, ..., x_4) = [3, 4, 2, 5, 1]

  * Gap-2 chain through position 0 is `x_0, x_2, x_4 = 3, 2, 1`.
  * HYPOTHESIS HOLDS.  At the first insertion step (`i = m = 2`) we compare `x_2 = 2` with
    `x_0 = 3`; since `2 < 3` they exchange places.  The state after that single step is
    `[2, 4, 3, 5, 1]` -- position 0 now holds `x_m = 2` and position `m` holds `x_0 = 3`, so
    the exchange is exhibited by RUNNING the algorithm, not asserted
    (`first_step_really_exchanges`).
  * CONCLUSION FAILS.  The fully m-sorted array is `B = [1, 4, 2, 5, 3]`, so `B[0] = 1`
    whereas the claim asserts `B[0] = x_m = 2`.
  * DIAGNOSIS CONFIRMED.  `B[0] = 1 = x_4 = x_{2m} = min{3, 2, 1}`, the minimum of the whole
    gap-`m` chain -- exactly the failure mode named in the record.

The record's smallest (degenerate, `m = 1`) witness `N = 3`, `A = [3, 2, 1]` is certified too
(`degenerate_witness_m_eq_1`): `x_1 = 2 < 3 = x_0` so they swap on the first step, yet
`B = [1, 2, 3]` and `B[0] = 1 = x_2 != x_1 = 2`.

## What is certified by kernel `decide` in this file

  (1) `msort_of_witness` / `two_implementations_agree_on_witness`: `B = [1,4,2,5,3]`, computed
      by BOTH implementations the record used -- the gap-`m` insertion pass (`msortIns`) and
      independent per-residue-chain sorting (`msortChain`) -- which agree.
  (2) `first_step_really_exchanges`: the hypothesis is verified by executing exactly one
      insertion step and observing the exchange of positions 0 and `m`.
  (3) `hypothesis_holds_conclusion_fails` / `claim_false_on_witness`: hypothesis true,
      conclusion false, on the same instance.
  (4) `witness_diagnosis`: `B[0]` equals the chain minimum and equals `x_{2m}`.
  (5) `gap_sort_first_entry_claim_is_false` / `stated_implication_fails`: the universally
      quantified claim is false, by instantiation at the witness.
  (6) Exhaustive sweeps (all 120 permutations of `[1..5]`; all 24 permutations of `[1..4]`;
      all 243 length-5 arrays over `{1,2,3}`, i.e. WITH ties) reproducing, at kernel scale,
      the record's "0 exceptions out of 36243 counterexamples" control.

## Modelling judgments (a reviewer should check each)

  1. ARRAYS are `List Nat`, indexed by `nth` (out-of-range reads return `0`) and updated by
     `setAt` (out-of-range writes are no-ops).  Every index used in every theorem is in
     range, so these conventions are never load-bearing.
  2. VALUES are `Nat`, so "<" is exact integer comparison; no floats, no sentinels.
  3. `msortIns m a` is the textbook gap-`m` insertion pass, transcribed literally:
     `for i in m..N-1: key := a[i]; j := i; while j >= m and a[j-m] > key: a[j] := a[j-m];
     j := j-m; a[j] := key`.  The comparison is STRICT (`a[j-m] > key`), i.e. the pass is
     stable -- the standard convention, and the one under which the source's "they exchange
     places" is a real exchange rather than a tie-break.
  4. Both loops are made structurally recursive by an explicit FUEL argument (Lean core has no
     `while`).  The fuel supplied is provably generous: the outer loop gets `N` (it runs
     `N - m` times) and each inner shift gets `i + 1` (`j` strictly decreases by `m >= 1` per
     iteration, starting from `j = i`).  Running out of fuel would only ever TRUNCATE the
     computation; that it does not is witnessed independently by (1) -- the fuel-driven
     implementation agrees with the fuel-independent chain-sorting implementation on the
     witness and on every array of all three exhaustive sweeps.
  5. `msortChain m a` is the SECOND, independent implementation: for each residue
     `r = m-1, ..., 0` it gathers the values at indices `r, r+m, r+2m, ...`, sorts them with a
     plain insertion sort, and scatters them back.  It shares no code with `msortIns` beyond
     `nth`/`setAt`.  This mirrors the record's own two-implementation cross-check.
  6. The HYPOTHESIS is modelled in the strong, literal form `hypExchangeB`: run one insertion
     step and demand that positions 0 and `m` have genuinely exchanged distinct values.  The
     source's own paraphrase `x_m < x_0` is modelled separately (`hypLessB`) and the two are
     certified EQUIVALENT across every sweep (`nc_hypothesis_paraphrase_is_faithful`), so the
     refutation cannot be blamed on reading the hypothesis too weakly.
  7. The CONCLUSION is modelled as literally as possible: `B[0] = x_m`, where `x_m` is the
     value at index `m` of the ORIGINAL array `A` (not of `B`).
  8. `m = 0` is degenerate (`msortIns 0` is the identity); it is excluded from every sweep,
     and the witness uses `m = 2 >= 2`, a genuine gap, so nothing here depends on the
     degenerate case.  The record's `m = 1` witness (plain insertion sort) is included only as
     a secondary confirmation.
  9. HUMAN-ADJUDICATION STEP that Lean cannot discharge: Lean certifies the crisp proposition
     "there is an array and a gap at which the first insertion step exchanges positions 0 and
     m yet the m-sorted array does not begin with x_m".  That this is what the source sentence
     asserts -- rather than, say, loose shorthand for "B[0] is the smaller of x_0 and x_m at
     that moment" -- is a reading of English, recorded in `caveats`, not a kernel fact.

## Negative controls (each would FAIL to compile if the modelling were vacuous or rigged)

  * `nc_claim_is_not_always_false`: on `[2,3,1,4,5]` with `m = 2` the hypothesis HOLDS and the
    claim's conclusion HOLDS (`B[0] = 1 = x_2`).  The encoding therefore does not make the
    conclusion unsatisfiable; the witness is a genuine, localised failure.
  * `nc_hypothesis_discriminates`: the hypothesis is neither constantly true nor constantly
    false -- it fails on the already-2-sorted `[1,4,2,5,3]` and holds on the witness; over the
    120 permutations of `[1..5]` it holds for some and fails for others.
  * `nc_msort_is_not_degenerate`: `msortIns 2` is neither the identity nor a full sort on the
    witness, and `msortIns 0` IS the identity (so the degenerate gap is visibly excluded).
  * `nc_msort_permutes_its_input`: over all 120 permutations, the multiset of entries is
    preserved -- the model does not lose or invent values.
  * `nc_msort_really_m_sorts`: over all 120 permutations, every residue chain of the output is
    sorted (so the pass does what "m-sort" means), while the output is NOT fully sorted for
    some input (so it is not secretly a plain sort).
  * `nc_two_implementations_agree_*`: the fuelled insertion pass and the chain sorter agree on
    all 120 permutations of `[1..5]` (`m = 2`), on all 24 permutations of `[1..4]` for every
    `m` in `{1,2,3}`, and on all 243 length-5 arrays over `{1,2,3}` (ties included).
  * `nc_repaired_*`: the REPAIRED statement `B[0] = min(chain through 0)` holds with ZERO
    exceptions on all 120 permutations for every `m` in `{1,2,3,4}`, and on all 243 arrays
    with ties -- reproducing the record's 0-exceptions diagnosis control at kernel scale.
  * `nc_claim_fails_and_holds_within_one_sweep`: within the SAME sweep the flagged claim has
    genuine violations AND genuine satisfied instances, so neither sweep above is true merely
    because its hypotheses are never met.

Everything is `Bool`-valued and structurally recursive over lists of length <= 5; kernel
evaluation is a few hundred thousand `Nat`/list steps.
-/

set_option maxRecDepth 20000
set_option maxHeartbeats 4000000

namespace QEDCert.SE_cs_81391__q

/-! ## 0. Tiny hand-rolled list plumbing (no `List` API is relied on) -/

def lapp {α : Type} : List α → List α → List α
  | [],      ys => ys
  | x :: xs, ys => x :: lapp xs ys

def lmap {α β : Type} : List α → (α → β) → List β
  | [],      _ => []
  | x :: xs, f => f x :: lmap xs f

def joinLL {α : Type} : List (List α) → List α
  | []      => []
  | l :: ls => lapp l (joinLL ls)

def lall {α : Type} : List α → (α → Bool) → Bool
  | [],      _ => true
  | x :: xs, p => p x && lall xs p

def lany {α : Type} : List α → (α → Bool) → Bool
  | [],      _ => false
  | x :: xs, p => p x || lany xs p

def llen {α : Type} : List α → Nat
  | []      => 0
  | _ :: xs => Nat.succ (llen xs)

/-- `[0, 1, ..., k-1]`. -/
def upto : Nat → List Nat
  | 0     => []
  | k + 1 => lapp (upto k) [k]

/-- Array read; out-of-range reads return `0` (never exercised below). -/
def nth : List Nat → Nat → Nat
  | [],      _     => 0
  | x :: _,  0     => x
  | _ :: xs, k + 1 => nth xs k

/-- Array write; out-of-range writes are no-ops (never exercised below). -/
def setAt : List Nat → Nat → Nat → List Nat
  | [],      _,     _ => []
  | _ :: xs, 0,     v => v :: xs
  | x :: xs, k + 1, v => x :: setAt xs k v

/-- Structural equality of `Nat` arrays. -/
def eqL : List Nat → List Nat → Bool
  | [],      []      => true
  | x :: xs, y :: ys => Nat.beq x y && eqL xs ys
  | _,       _       => false

/-- Is the list nondecreasing? -/
def isSortedB : List Nat → Bool
  | []           => true
  | [_]          => true
  | x :: y :: r  => Nat.ble x y && isSortedB (y :: r)

/-! ## 1. Plain insertion sort (used only by the CHAIN implementation and by the controls) -/

def insSorted (x : Nat) : List Nat → List Nat
  | []      => [x]
  | a :: as =>
      match Nat.ble x a with
      | true  => x :: a :: as
      | false => a :: insSorted x as

def sortNat : List Nat → List Nat
  | []      => []
  | a :: as => insSorted a (sortNat as)

/-! ## 2. Implementation A -- the gap-`m` insertion pass, transcribed literally

    for i := m to N-1:
        key := a[i]; j := i
        while j >= m and a[j-m] > key:  a[j] := a[j-m];  j := j - m
        a[j] := key
-/

/-- The inner `while` loop: shift chain elements right while they exceed `key`, then drop
`key` into the hole.  `fuel` bounds the number of shifts (`j` decreases by `m` each time). -/
def shiftIns (m key : Nat) : List Nat → Nat → Nat → List Nat
  | a, j, 0          => setAt a j key
  | a, j, Nat.succ f =>
      match Nat.ble m j && Nat.blt key (nth a (j - m)) with
      | true  => shiftIns m key (setAt a j (nth a (j - m))) (j - m) f
      | false => setAt a j key

/-- The outer `for` loop over `i = m, m+1, ...`, bounded by `fuel` iterations. -/
def msLoop (m n : Nat) : List Nat → Nat → Nat → List Nat
  | a, _, 0          => a
  | a, i, Nat.succ f =>
      match Nat.blt i n with
      | false => a
      | true  => msLoop m n (shiftIns m (nth a i) a i (Nat.succ i)) (Nat.succ i) f

/-- **m-sort, implementation A**: the full gap-`m` insertion pass. -/
def msortIns (m : Nat) (a : List Nat) : List Nat :=
  msLoop m (llen a) a m (llen a)

/-- The array after exactly the FIRST insertion step (`i = m`) -- used to exhibit the
claim's hypothesis by execution. -/
def msortFirstStep (m : Nat) (a : List Nat) : List Nat :=
  msLoop m (llen a) a m 1

/-! ## 3. Implementation B -- sort each residue chain in place (independent of A) -/

/-- Indices of the gap-`m` chain through `r`: `r, r+m, r+2m, ...` below `n`. -/
def chainIdx (n m : Nat) : Nat → Nat → List Nat
  | _, 0          => []
  | r, Nat.succ f =>
      match Nat.blt r n with
      | false => []
      | true  => r :: chainIdx n m (r + m) f

def gather (a : List Nat) : List Nat → List Nat
  | []      => []
  | i :: is => nth a i :: gather a is

def scatter : List Nat → List Nat → List Nat → List Nat
  | a, [],      _       => a
  | a, _ :: _,  []      => a
  | a, i :: is, v :: vs => scatter (setAt a i v) is vs

/-- Sort the single residue chain through `r`. -/
def sortChain (m n : Nat) (a : List Nat) (r : Nat) : List Nat :=
  scatter a (chainIdx n m r n) (sortNat (gather a (chainIdx n m r n)))

/-- Sort chains `r = k-1, k-2, ..., 0`. -/
def chainLoop (m n : Nat) : List Nat → Nat → List Nat
  | a, 0          => a
  | a, Nat.succ r => chainLoop m n (sortChain m n a r) r

/-- **m-sort, implementation B**: sort every gap-`m` residue chain in place. -/
def msortChain (m : Nat) (a : List Nat) : List Nat :=
  chainLoop m (llen a) a m

/-- Every residue chain of `a` is sorted. -/
def allChainsSortedB (m : Nat) (a : List Nat) : Bool :=
  lall (upto m) (fun r => isSortedB (gather a (chainIdx (llen a) m r (llen a))))

/-! ## 4. The claim, as decidable predicates -/

/-- **The claim's hypothesis, literal form**: at the first insertion step, positions `0` and
`m` genuinely exchange (distinct) values.  Verified by running the step. -/
def hypExchangeB (m : Nat) (a : List Nat) : Bool :=
  Nat.blt m (llen a)
    && Nat.beq (nth (msortFirstStep m a) 0) (nth a m)
    && Nat.beq (nth (msortFirstStep m a) m) (nth a 0)
    && !(Nat.beq (nth a 0) (nth a m))

/-- **The source's own paraphrase of the hypothesis**: `x_m < x_0`. -/
def hypLessB (m : Nat) (a : List Nat) : Bool :=
  Nat.blt m (llen a) && Nat.blt (nth a m) (nth a 0)

/-- **The claim's conclusion**: the m-sorted array begins with `x_m`. -/
def conclB (m : Nat) (a : List Nat) : Bool :=
  Nat.beq (nth (msortIns m a) 0) (nth a m)

/-- **The claim at one instance**: hypothesis implies conclusion. -/
def claimAt (m : Nat) (a : List Nat) : Bool :=
  match hypExchangeB m a with
  | false => true
  | true  => conclB m a

/-- The minimum of the gap-`m` chain through position `0`. -/
def chainMin (m : Nat) (a : List Nat) : Nat :=
  nth (sortNat (gather a (chainIdx (llen a) m 0 (llen a)))) 0

/-- **The REPAIRED statement** (the record's diagnosis): `B[0] = min{x_0, x_m, x_{2m}, ...}`.
Note it is unconditional -- no hypothesis needed. -/
def repairedAt (m : Nat) (a : List Nat) : Bool :=
  Nat.beq (nth (msortIns m a) 0) (chainMin m a)

/-! ## 5. The witness -/

/-- `A = [3, 4, 2, 5, 1]`, `N = 5`. -/
def A : List Nat := [3, 4, 2, 5, 1]

/-- The gap `m = 2` -- a genuine gap (`m >= 2`), not the degenerate insertion-sort case. -/
def mW : Nat := 2

/-- `B = m-sort(A)`. -/
def B : List Nat := msortIns mW A

/-! ## 6. The refutation -/

/-- `B = [1, 4, 2, 5, 3]`. -/
theorem msort_of_witness : eqL B [1, 4, 2, 5, 3] = true := by decide

/-- The two independent implementations agree on the witness (the record's cross-check). -/
theorem two_implementations_agree_on_witness :
    eqL (msortIns mW A) (msortChain mW A) = true := by decide

/-- **Hypothesis, by execution.**  After exactly one insertion step the array is
`[2, 4, 3, 5, 1]`: position `0` holds `x_m = 2` and position `m` holds `x_0 = 3`, so `x_0` and
`x_m` really did "switch place during the first pass"; the source's paraphrase `x_m < x_0`
holds as well. -/
theorem first_step_really_exchanges :
    (eqL (msortFirstStep mW A) [2, 4, 3, 5, 1]
      && hypExchangeB mW A
      && hypLessB mW A) = true := by decide

/-- **MAIN (instance form).**  The hypothesis holds and the conclusion fails: `B[0] = 1`,
whereas the claim asserts `B[0] = x_m = 2`. -/
theorem hypothesis_holds_conclusion_fails :
    (hypExchangeB mW A
      && !(conclB mW A)
      && Nat.beq (nth B 0) 1
      && Nat.beq (nth A mW) 2) = true := by decide

/-- The same, as the claim evaluating to `false` on the witness. -/
theorem claim_false_on_witness : claimAt mW A = false := by decide

/-- **The record's diagnosis, confirmed.**  `B[0]` is the minimum of the WHOLE gap-`m` chain
through position `0` (`{x_0, x_2, x_4} = {3, 2, 1}`), and here that minimum is `x_{2m} = x_4 =
1`, not `x_m`. -/
theorem witness_diagnosis :
    (Nat.beq (nth B 0) (chainMin mW A)
      && Nat.beq (chainMin mW A) 1
      && Nat.beq (nth A 4) 1
      && eqL (gather A (chainIdx (llen A) mW 0 (llen A))) [3, 2, 1]) = true := by decide

/-- **MAIN (universal form).**  The claim, universally quantified over gaps and arrays, is
false. -/
theorem gap_sort_first_entry_claim_is_false :
    ¬ ∀ (m : Nat) (a : List Nat), claimAt m a = true := by
  intro h
  exact absurd (h mW A) (by decide)

/-- **MAIN (implication form).**  Stated the way the source states it: "if `x_0` and `x_m`
exchange on the first step then `B[0] = x_m`". -/
theorem stated_implication_fails :
    ¬ ∀ (m : Nat) (a : List Nat),
        hypExchangeB m a = true → Nat.beq (nth (msortIns m a) 0) (nth a m) = true := by
  intro h
  exact absurd (h mW A (by decide)) (by decide)

/-- The record's smallest, degenerate witness (`m = 1`, plain insertion sort): `A' = [3,2,1]`,
the first step swaps `x_0` and `x_1`, yet `B' = [1,2,3]` and `B'[0] = 1 = x_2 != 2 = x_1`. -/
theorem degenerate_witness_m_eq_1 :
    (hypExchangeB 1 [3, 2, 1]
      && eqL (msortIns 1 [3, 2, 1]) [1, 2, 3]
      && eqL (msortChain 1 [3, 2, 1]) [1, 2, 3]
      && !(conclB 1 [3, 2, 1])
      && claimAt 1 [3, 2, 1] == false) = true := by decide

/-! ## 7. Exhaustive sweeps -/

/-- All insertions of `x` into `l`. -/
def insertsOf (x : Nat) : List Nat → List (List Nat)
  | []      => [[x]]
  | a :: as => (x :: a :: as) :: lmap (insertsOf x as) (fun t => a :: t)

/-- All permutations of a list. -/
def permsOf : List Nat → List (List Nat)
  | []      => [[]]
  | a :: as => joinLL (lmap (permsOf as) (fun p => insertsOf a p))

/-- All 120 permutations of `[1,2,3,4,5]`. -/
def perms5 : List (List Nat) := permsOf [1, 2, 3, 4, 5]

/-- All 24 permutations of `[1,2,3,4]`. -/
def perms4 : List (List Nat) := permsOf [1, 2, 3, 4]

/-- All length-`k` tuples over the alphabet `vals` (arrays WITH ties). -/
def tuples (vals : List Nat) : Nat → List (List Nat)
  | 0     => [[]]
  | k + 1 => joinLL (lmap (tuples vals k) (fun t => lmap vals (fun v => v :: t)))

/-- All 243 length-5 arrays over `{1,2,3}`. -/
def ties5 : List (List Nat) := tuples [1, 2, 3] 5

/-- A counterexample WITH repeated keys, used to show `ties5` is populated by relevant
arrays: `[3,1,2,1,1]` at `m = 2` m-sorts to `[1,1,2,1,3]`, so `B[0] = 1 != 2 = x_m`. -/
def Aties : List Nat := [3, 1, 2, 1, 1]

theorem nc_sweeps_are_populated :
    (Nat.beq (llen perms5) 120
      && Nat.beq (llen perms4) 24
      && Nat.beq (llen ties5) 243
      && lany perms5 (fun a => eqL a A)
      && lany ties5 (fun a => eqL a Aties)
      && claimAt mW Aties == false) = true := by decide

/-! ## 8. Negative controls -/

/-- **NEGATIVE CONTROL 1 (the headline control).**  The claim's conclusion is NOT always
false under this encoding: on `[2,3,1,4,5]` with `m = 2` the hypothesis holds (`1 < 2`) AND
the conclusion holds (`B = [1,3,2,4,5]`, `B[0] = 1 = x_2`).  If the modelling were rigged --
e.g. if `conclB` were secretly the negation of something -- this theorem would not compile,
and the witness would be an artefact rather than a genuine localised failure. -/
theorem nc_claim_is_not_always_false :
    (hypExchangeB 2 [2, 3, 1, 4, 5]
      && conclB 2 [2, 3, 1, 4, 5]
      && claimAt 2 [2, 3, 1, 4, 5]
      && eqL (msortIns 2 [2, 3, 1, 4, 5]) [1, 3, 2, 4, 5]) = true := by decide

/-- **NEGATIVE CONTROL 2.**  The hypothesis discriminates: it fails on the already-2-sorted
`[1,4,2,5,3]` (there `x_2 = 2 > 1 = x_0`, no exchange) and holds on the witness; and both
polarities occur across the 120-permutation sweep. -/
theorem nc_hypothesis_discriminates :
    (hypExchangeB mW [1, 4, 2, 5, 3] == false
      && hypExchangeB mW A
      && lany perms5 (fun a => hypExchangeB mW a)
      && lany perms5 (fun a => !(hypExchangeB mW a))) = true := by decide

/-- **NEGATIVE CONTROL 3.**  The literal hypothesis and the source's paraphrase `x_m < x_0`
agree on every one of the 120 permutations, so the refutation does not turn on reading
"switches place with" too weakly or too strongly. -/
theorem nc_hypothesis_paraphrase_is_faithful :
    lall perms5 (fun a => hypExchangeB mW a == hypLessB mW a) = true := by decide

/-- **NEGATIVE CONTROL 4.**  `msortIns 2` is neither the identity nor a full sort on the
witness -- it is a genuine partial (gap-2) sort; and `msortIns 0` IS the identity, which is
why the degenerate gap is excluded from every sweep. -/
theorem nc_msort_is_not_degenerate :
    (eqL (msortIns mW A) A == false
      && isSortedB (msortIns mW A) == false
      && isSortedB A == false
      && eqL (msortIns 0 A) A) = true := by decide

/-- **NEGATIVE CONTROL 5.**  The pass really m-sorts: over all 120 permutations every residue
chain of the output is sorted... -/
theorem nc_msort_really_m_sorts :
    lall perms5 (fun a => allChainsSortedB mW (msortIns mW a)) = true := by decide

/-- ...while the output is NOT fully sorted for some input, so `msortIns` is not secretly an
ordinary sort (which would make `B[0]` the global minimum for trivial reasons). -/
theorem nc_msort_is_not_a_full_sort :
    lany perms5 (fun a => !(isSortedB (msortIns mW a))) = true := by decide

/-- **NEGATIVE CONTROL 6.**  The pass loses and invents nothing: over all 120 permutations the
sorted output equals the sorted input. -/
theorem nc_msort_permutes_its_input :
    lall perms5 (fun a => eqL (sortNat (msortIns mW a)) (sortNat a)) = true := by decide

/-- **NEGATIVE CONTROL 7a.**  Implementations A and B agree on all 120 permutations
(`m = 2`) -- so the fuel bounds in A never truncate and the transcription is faithful. -/
theorem nc_two_implementations_agree_perms5 :
    lall perms5 (fun a => eqL (msortIns mW a) (msortChain mW a)) = true := by decide

/-- **NEGATIVE CONTROL 7b.**  ...and on all 24 permutations of `[1..4]` for every gap
`m` in `{1,2,3}`. -/
theorem nc_two_implementations_agree_perms4_all_gaps :
    lall perms4 (fun a =>
        lall ([1, 2, 3] : List Nat) (fun m => eqL (msortIns m a) (msortChain m a)))
      = true := by decide

/-- **NEGATIVE CONTROL 7c.**  ...and on all 243 length-5 arrays over `{1,2,3}`, i.e. WITH
repeated keys, where stability could have separated the two implementations. -/
theorem nc_two_implementations_agree_with_ties :
    lall ties5 (fun a => eqL (msortIns mW a) (msortChain mW a)) = true := by decide

/-- **NEGATIVE CONTROL 8a (the record's diagnosis control).**  The REPAIRED statement
`B[0] = min(gap-m chain through 0)` holds with ZERO exceptions on all 120 permutations of
`[1..5]`, for each gap `m` in `{1,2,3,4}` (split into one theorem per gap purely to keep each
kernel `decide` small).  This is the kernel-scale reproduction of the record's
"0 exceptions over 36243 counterexamples" check, and it pins the failure on the claim's
identification of `B[0]` with `x_m` rather than on a broken model of m-sorting. -/
theorem nc_repaired_holds_on_perms5_gap1 :
    lall perms5 (fun a => repairedAt 1 a) = true := by decide

theorem nc_repaired_holds_on_perms5_gap2 :
    lall perms5 (fun a => repairedAt 2 a) = true := by decide

theorem nc_repaired_holds_on_perms5_gap3 :
    lall perms5 (fun a => repairedAt 3 a) = true := by decide

theorem nc_repaired_holds_on_perms5_gap4 :
    lall perms5 (fun a => repairedAt 4 a) = true := by decide

/-- **NEGATIVE CONTROL 8b.**  The same with ties: zero exceptions on all 243 length-5 arrays
over `{1,2,3}`. -/
theorem nc_repaired_holds_with_ties :
    lall ties5 (fun a => repairedAt mW a) = true := by decide

/-- **NEGATIVE CONTROL 9.**  Within the SAME sweep the flagged claim genuinely FAILS on some
arrays and genuinely HOLDS (with its hypothesis satisfied) on others.  Without this, the
sweeps above could be true merely because the hypothesis is never met, or the refutation could
be an artefact of a hypothesis that is met everywhere. -/
theorem nc_claim_fails_and_holds_within_one_sweep :
    (lany perms5 (fun a => !(claimAt mW a))
      && lany perms5 (fun a => hypExchangeB mW a && conclB mW a)
      && lany ties5 (fun a => !(claimAt mW a))) = true := by decide

/-! ## 9. Bundled certificate -/

/-- **Bundled certificate.**  On `A = [3,4,2,5,1]` with `m = 2`: one insertion step exchanges
`x_0 = 3` and `x_m = 2` (so the claim's hypothesis holds); the fully m-sorted array is
`B = [1,4,2,5,3]`, computed identically by two independent implementations; `B[0] = 1 != 2 =
x_m`, so the claim's conclusion fails; and `B[0]` is the minimum `1 = x_{2m}` of the whole
gap-`m` chain `{3,2,1}`, exactly as the record diagnoses. -/
theorem certificate :
    eqL (msortFirstStep mW A) [2, 4, 3, 5, 1] = true
      ∧ hypExchangeB mW A = true
      ∧ hypLessB mW A = true
      ∧ eqL B [1, 4, 2, 5, 3] = true
      ∧ eqL (msortIns mW A) (msortChain mW A) = true
      ∧ conclB mW A = false
      ∧ claimAt mW A = false
      ∧ Nat.beq (nth B 0) (chainMin mW A) = true :=
  ⟨by decide, by decide, by decide, by decide, by decide, by decide, by decide, by decide⟩

end QEDCert.SE_cs_81391__q

\end{Vlean}
\subsection*{Proof-safety information}
The automated scan found no unfinished proofs, new axiom declarations, unsafe declarations, external code hooks, or uses of Lean's compiled decision procedure. This scan concerns proof construction only; it does not establish that the theorem matches the source claim.
\textit{Technical scan counts:} native\_decide=0, sorry=0, axiom\_decl=0, unsafe=0, extern=0. \textit{Technical tier:} KERNEL-ONLY.
\textbf{A YES means:} Lean represents the same objects and assumptions, uses the same proposed counterexample, and the final theorem statement asserts the same contradiction. Successful compilation by itself is not enough.
\clearpage
\section*{Human sign-off}
Complete this block from your own review of Questions 1--3. Do not use the automated appendix below as an answer key.\par
\textbf{Reviewer name:} \rule{9cm}{0.4pt}\par
\textbf{Date:} \rule{4cm}{0.4pt}\par
\medskip\textbf{1. Does the proof verification document accurately represent the source and challenged claim?}\par
\answerchoices\par
\decisionexplanation
\medskip\textbf{2. Does the proposed counterexample satisfy the assumptions and make that claim false?}\par
\answerchoices\par
\decisionexplanation
\medskip\textbf{3. Does the final Lean theorem statement represent and verify that same counterexample?}\par
\answerchoices\par
\decisionexplanation
\medskip\textbf{Overall decision:}\par
\checkbox\ VALIDATED (all three answers are YES)\par
\checkbox\ REJECTED (at least one answer is NO)\par
\checkbox\ NEEDS ANOTHER REVIEWER (at least one CANNOT DETERMINE and no NO)\par
\medskip\textbf{Additional comments.} Use this space for any extra source, counterexample, calculation, or Lean details that support your decisions.\par\rule{\linewidth}{0.4pt}\par\vspace{7pt}\rule{\linewidth}{0.4pt}
\clearpage
\section*{Appendix --- Prior automated review (background only)}
\textbf{These results are not an answer key.} They come from the earlier Opus/Fable verification campaign and are printed only so an auditor can see what was previously checked. The human decisions above must be made independently. Raw status codes are retained in small type for traceability.
\textbf{Overall automated disposition:} Cleared by automated screening; an independent search also found a counterexample.
\par\textit{Audit code:} \nolinkurl{PASS_CORROBORATED}
\subsection*{Independent search with no access to the stored example}
\textbf{Result:} An independent search found a counterexample without seeing the stored example.
\par\textit{Audit code:} \nolinkurl{WITNESS_FOUND}
\par\textbf{Checker's explanation:}
\begin{Vsrc}
The claim's parenthetical is self-undermining: it concedes that after x_0 is displaced to position m it may 'further switch place with x_{2m}', but that is precisely the exchange that hands x_{2m} a path to position 0. In the gap-m insertion pass, index 2m is processed after index m, and if x_{2m} is smaller than the current occupant of position 0 the insertion loop shifts it all the way to position 0. So position 0 ends up holding the minimum of the entire chain, not x_m. (Shortened here; the complete note remains in the machine-check JSON.)
\end{Vsrc}
\par\textbf{Cases checked:}
\begin{Vsrc}
Deterministic exhaustive sweep, stdlib only, integers only. For every length n in 3..7 and every gap m in 1..n-1: every value tuple in {0,...,n-1}^n (all n**n arrays, ties included) plus every permutation of {0,...,n-1} (tie-free sub-sweep reported separately). 5,222,264 arrays swept in total. (Shortened here; the complete note remains in the machine-check JSON.)
\end{Vsrc}
\subsection*{Independent re-check of the proposed counterexample}
\textbf{Result:} A separate checker recomputed the proposed counterexample and confirmed it.
\par\textit{Audit code:} \nolinkurl{CONFIRMED}
\par\textbf{Checker's note:}
\begin{Vsrc}
CONFIRMED. My independent recomputation reproduces the recorded witness exactly: for A = [3,4,2,5,1] with m = 2 the hypothesis of the quoted claim genuinely holds (x_m = 2 < 3 = x_0, so at the first insertion step i = m the entries at positions 0 and m exchange), yet m-sort(A) = [1,4,2,5,3], so B[0] = 1 while the claim asserts B[0] = x_m = 2. (Shortened here; the complete note remains in the machine-check JSON.)
\end{Vsrc}
\subsection*{Independent review of the source-to-Lean correspondence}
\textbf{Result:} A separate reviewer judged that Lean proves the same failure using the same example.
\par\textit{Audit code:} \nolinkurl{MATCHES}
\par\textbf{Theorem summary from that reviewer:}
\begin{Vsrc}
The Lean file defines m-sorting as the textbook gap-m insertion pass (`msortIns`, fuel-bounded structural recursion) plus an independent per-residue-chain sorter (`msortChain`), and encodes the flagged inference as `hypExchangeB m a` (running exactly one insertion step at i = m leaves position 0 holding the original a[m] and position m holding the original a[0], with a[0]!= a[m]) implying `conclB m a` ((msortIns m a)[0] = a[m], where a[m] = x_m is read off the ORIGINAL array). (Shortened here; the complete note remains in the machine-check JSON.)
\end{Vsrc}
\par\textbf{Difference noted by that reviewer:}
\begin{Vsrc}
None that break the three links. Verified independently (Python reimplementation of both a gap-m insertion pass and per-chain sorting): (1) DEFINITIONS -- msortIns is the literal gap-m insertion pass with strict comparison; it agrees with the independent chain sorter on the witness, on all 120 permutations of [1..5], on all 24 permutations of [1..4] for m in {1,2,3}, and on all 243 length-5 arrays over {1,2,3} (ties), so the fuel bounds never truncate; (Shortened here; the complete note remains in the machine-check JSON.)
\end{Vsrc}
\subsection*{Follow-up review}
No follow-up review was needed or recorded.
\textit{This appendix is only background. You do not need any separate machine-output files to answer the three questions.}
\end{document}